% thesis/figures/fig3_calib_pipeline.tex
% Figure: 行为引导离线校准流程 (Ch3 §3.3 对应)
% 两条路径共用 KL 散度作为校准目标

\begin{figure}[!htbp]
  \centering
  \begin{tikzpicture}[
    every node/.style={font=\small},
    data/.style={rectangle, rounded corners=2pt, draw=gray!70, thick,
                 fill=gray!10, minimum width=2.4cm, minimum height=0.9cm, align=center},
    fp/.style={rectangle, rounded corners=2pt, draw=blue!55!black, thick,
               fill=blue!10, minimum width=2.6cm, minimum height=0.9cm, align=center},
    int8/.style={rectangle, rounded corners=3pt, draw=green!50!black, very thick,
                 fill=green!12, minimum width=4.2cm, minimum height=1.5cm, align=center},
    int4/.style={rectangle, rounded corners=3pt, draw=orange!80!black, very thick,
                 fill=orange!15, minimum width=4.2cm, minimum height=1.5cm, align=center},
    art/.style={rectangle, rounded corners=2pt, draw=purple!60, thick,
                fill=purple!8, minimum width=3.2cm, minimum height=1.0cm, align=center},
    goal/.style={rectangle, rounded corners=3pt, draw=red!60, thick,
                 fill=red!8, minimum width=8cm, minimum height=0.8cm, align=center},
    arr/.style={-stealth, thick},
    sharr/.style={-stealth, thick, dashed, gray!70},
  ]

    % ──── Step 1: 输入（校准数据 + FP16 前向）────
    \node[data] (data) at (-5.2, 3.2) {WikiText-2\\ 校准数据};
    \node[fp]   (fwd)  at (-1.0, 3.2) {FP16 前向\\ 提取 $(Q, K, V)$};
    \draw[arr] (data.east) -- (fwd.west);

    % ──── Step 2: 校准目标 (KL 散度最小化) ────
    \node[goal] (goal) at (-1.0, 2.0) {
      \textbf{校准目标：}\ $D_{\mathrm{KL}}\!\bigl(\bm{a}_{\mathrm{ref}} \,\|\, \bm{a}_{\mathrm{quant}}\bigr)$ 最小化\
      \footnotesize (注意力分布偏移度)
    };
    \draw[arr] (fwd.south) -- (goal.north);

    % ──── Step 3: 两路搜索（左 INT8 / 右 INT4-RoleAlign）────
    \node[int8] (int8) at (-5.3, 0.0) {
      \textbf{INT8 对称路径}\\[3pt]
      \footnotesize 搜索 $(p_c, g)$ 网格\\
      \footnotesize 对称 per-group 量化
    };
    \node[int4] (int4) at (3.2, 0.0) {
      \textbf{INT4-RoleAlign 非对称路径}\\[3pt]
      \footnotesize 搜索 $(p_K, p_V)$\\
      \footnotesize per-channel K + per-token V
    };
    \draw[arr] (goal.south) -| (int8.north);
    \draw[arr] (goal.south) -| (int4.north);

    % ──── Step 4: 输出（校准产物：参数类型而非文件名）────
    \node[art] (art_int8) at (-5.3, -1.8) {
      \textbf{INT8 校准产物}\\
      \footnotesize 逐层 per-group 静态 Scale
    };
    \node[art] (art_int4) at (3.2, -1.8) {
      \textbf{INT4-RoleAlign 校准产物}\\
      \footnotesize 逐层 $(p_K, p_V)$ 参数
    };
    \draw[arr] (int8.south) -- (art_int8.north);
    \draw[arr] (int4.south) -- (art_int4.north);

    % ──── 连线：两路共享同一校准目标（纯箭头, 无冗余 label）────
    \draw[sharr] (int8.east) -- (int4.west);

  \end{tikzpicture}
  \caption{行为引导离线校准流程。两条路径（INT8 对称 / INT4-RoleAlign 非对称）
  共用 KL 散度作为校准目标——对注意力分布偏移度取最小；搜索空间的差异
  （$(p_c, g)$ 对称 per-group \textit{vs.} $(p_K, p_V)$ 非对称 per-channel/per-token）
  来自量化格式本身，而非目标函数。校准产物为推理时只读的离线参数集。}
  \label{fig:calib-pipeline}
\end{figure}
